\documentclass[11pt, a4paper]{article}

% --- UNIVERSAL PREAMBLE BLOCK ---
\usepackage[a4paper, top=2.2cm, bottom=2.2cm, left=2.2cm, right=2.2cm]{geometry}
\usepackage{fontspec}

\usepackage[italian, bidi=basic, provide=*]{babel}

\babelprovide[import, onchar=ids fonts]{italian}
\babelprovide[import, onchar=ids fonts]{english}

% Set default/Latin font to Sans Serif in the main (rm) slot
\babelfont{rm}{Noto Sans}

% Add because main language is not English
\usepackage{enumitem}
% Configurazione liste per renderle più eleganti
\setlist[itemize]{label=\textcolor{primary}{\textbullet}, leftmargin=*, itemsep=2pt, parsep=0pt}

% --- CUSTOM STYLING PACKAGES ---
\usepackage{xcolor}
\usepackage{titlesec}
\usepackage{tabularx}
\usepackage{parskip}

% I due colori sobri ed eleganti scelti
\definecolor{primary}{HTML}{1A4B84} % Blu Navy elegante
\definecolor{textdark}{HTML}{222222} % Antracite/Grigio scuro per il testo

% Imposta il colore base del testo
\color{textdark}

% Formattazione Titoli di Sezione
\titleformat{\section}
  {\Large\bfseries\color{primary}} % Stile
  {} % Nessun numero
  {0em} % Distanza dal margine
  {} % Codice prima del titolo
  [\vspace{-0.3em}\textcolor{primary}{\rule{\textwidth}{1pt}}] % Linea orizzontale sotto il titolo
\titlespacing*{\section}{0pt}{3.5ex}{1.5ex}

% Macro personalizzata per le voci di Esperienza e Istruzione
\newcommand{\resumeItem}[4]{
  \vspace{1.5mm}
  \noindent
  \begin{tabular*}{\textwidth}{@{}l@{\extracolsep{\fill}}r@{}}
    \textbf{#1} & \textbf{\textcolor{primary}{#2}} \\
    \textit{#3} & \textit{#4} \\
  \end{tabular*}
  \vspace{1mm}
}

% Hyperref (DEVE ESSERE L'ULTIMO)
\usepackage[colorlinks=true, urlcolor=primary, linkcolor=primary]{hyperref}

\begin{document}

% --- INTESTAZIONE CON SEGNAPOSTO FOTO ---
\noindent
\begin{minipage}[c]{0.75\textwidth}
    {\Huge \textbf{\textcolor{primary}{GIACOMO LOMBARDO}}}\\[0.5em]
    {\Large \textbf{AI / ML Engineer}}\\[1em]
    \small
    \href{mailto:giacomolombardo67@gmail.com}{giacomolombardo67@gmail.com} \quad|\quad 
    +39 3248703757 \\[0.3em]
    Via G. De Stefani N°4, 91029 Santa Ninfa (TP) \\[0.3em]
    Nato il 03/01/2003 \quad\\[0.3em]
    \href{https://www.linkedin.com/in/giacomo-lombardo-a944b8315/}{LinkedIn: https://www.linkedin.com/in/giacomo-lombardo-a944b8315/}\\
    \href{https://github.com/Giacomo0303}{GitHub: https://github.com/Giacomo0303}
\end{minipage}%
\begin{minipage}[c]{0.25\textwidth}
    \begin{flushright}
        % ISTRUZIONI PER LA FOTO REALE:
        % Quando compilerai il file sul tuo PC, cancella il comando \framebox qui sotto
        % e usa: \includegraphics[width=3cm]{nome_della_tua_foto.jpg}
        % (Ricordati di aggiungere \usepackage{graphicx} nel preambolo in alto!)
        \framebox{\parbox[c][4cm][c]{3cm}{\centering \textcolor{gray}{\small [Inserisci\\[0.2em]Foto Qui]}}}
    \end{flushright}
\end{minipage}
\vspace{1.5em}

% --- PROFILO ---
\section{Profilo Professionale}
Ingegnere Informatico (Laurea Triennale) attualmente in fase di specializzazione con un corso di Laurea Magistrale. Nutro un forte interesse per i campi dell'Intelligenza Artificiale, con focus specifico su \textbf{Computer Vision} ed \textbf{Elaborazione del Linguaggio Naturale (NLP)}. Dotato di una solida base algoritmica e architetturale, cerco una posizione come \textbf{AI/ML Engineer} per contribuire allo sviluppo di soluzioni all'avanguardia.

% --- ESPERIENZA ---
\section{Esperienza Professionale}

\resumeItem{Sviluppatore AI / Tirocinio Engineering}{Luglio 2023 -- Ottobre 2023}{ITALTEL S.p.A.}{Palermo (PA)}
Tirocinio curriculare tecnico (75 ore).
\begin{itemize}
    \item \textbf{Progetto principale:} Progettazione e sviluppo di un software di Computer Vision dedicato al riconoscimento automatico e al conteggio di etichette in input fotografici.
    \item \textbf{Stack e Competenze:} Sviluppo e addestramento di modelli di Machine Learning utilizzando \textbf{TensorFlow} e le API \textbf{Keras}.
    \item \textbf{Architetture:} Progettazione, implementazione e addestramento di \textbf{Reti Neurali Convoluzionali (CNN)} per l'analisi e classificazione delle immagini.
\end{itemize}

% --- ISTRUZIONE ---
\section{Istruzione e Formazione}

\resumeItem{Laurea Magistrale in Ingegneria Informatica (LM-32)}{2024 -- In corso}{Università degli Studi di Palermo}{Palermo}
\vspace{-1.5em}
\begin{itemize}
    \item Tesi basata su compressione di modelli Transformer attraverso tecniche di pruning e ricerca architetturale.
\end{itemize}

\resumeItem{Laurea Triennale in Ingegneria Informatica (L-8)}{2021 -- Luglio 2024}{Università degli Studi di Palermo}{Palermo}
\vspace{-1.5em}
\begin{itemize}
    \item Valutazione finale: \textbf{110/110 con lode}
\end{itemize}

\resumeItem{Diploma di Maturità Scientifica (Scienze Applicate)}{2016 -- 2021}{Liceo Scientifico Michele Cipolla}{Castelvetrano (TP)}
\vspace{-1.5em}
\begin{itemize}
    \item Valutazione finale: \textbf{100/100 con lode}
\end{itemize}

% --- COMPETENZE ---
\section{Competenze Tecniche}
\begin{itemize}
    \item \textbf{Intelligenza Artificiale:} Computer Vision, Natural Language Processing (NLP), Deep Learning, ML.
    \item \textbf{Framework e Librerie AI:} PyTorch, TensorFlow, Keras, Scikit-learn, Pandas, NumPy.
    \item \textbf{Linguaggi di Programmazione:} Python (Focus AI \& Data Science), Java, C, JavaScript.
    \item \textbf{Database, Dati \& Sistemi:} Database a grafo (Neo4j), Database Relazionali (SQL), JSON, XML, Linux, Windows.
    \item \textbf{Strumenti \& IDE:} PyCharm, Git.
\end{itemize}

% --- CERTIFICAZIONI E LINGUE ---
\section{Certificazioni e Lingue}
\begin{itemize}
    \item \textbf{Lingua Inglese:} Livello B2.
    \item \textbf{Lingua Italiana:} Madrelingua.
\end{itemize}

\vfill
\begin{center}
\scriptsize{\textit{Autorizzo il trattamento dei miei dati personali ai sensi del Decreto Legislativo 30 giugno 2003, n. 196 e del GDPR (Regolamento UE 2016/679).}}
\end{center}

\end{document}